\clearpage
\VerdierPhysicalPageAnchor{0240}%
\begingroup\small\setlength{\parskip}{0pt}\setlength{\emergencystretch}{1.5em}
\setlength{\abovedisplayskip}{4pt plus 2pt minus 2pt}%
\setlength{\belowdisplayskip}{4pt plus 2pt minus 2pt}%

\VerdierSemanticAnchor{ch4d-4-4-5}
\noindent\textbf{4.4.5.}\quad
データは (4.4.4) と同じであるが、さらに $D(\mathcal A)$ の対象 $X$ を
与える。$\bigl(X(p,q),\delta(p,q,r)\bigr)$ を、複体 $X$ に付随する、
$D(\mathcal A)$ に値をとる第二スペクトル対象とする (4.3.1)。関手 $F$ は
このスペクトル対象を $\mathcal B$ に値をとるスペクトル対象へ送る。
これにより、後者が圏 $D(\mathcal A)$ の中を動くとき、$F$ と引数 $X$ に
関手的なスペクトル系列が得られる (4.3.3)。このスペクトル系列を
\emph{複体 $X$ に関する関手 $F$ の第二スペクトル系列}という。この
スペクトル系列には常に（第 II 章、4.3.3.2）で導入した添字付けを用いる。

\VerdierSemanticAnchor{ch4d-4-4-6}
\medskip\noindent\textbf{命題 4.4.6.}\quad

\VerdierSemanticAnchor{ch4d-4-4-6-a}
\noindent\textit{a) 複体 $X$ に関する $F$ の第二スペクトル系列の項
$\mathrm{II}_2^{p,q}$ は次で与えられる：}
\[
\mathrm{II}_2^{p,q}=F^p\bigl(H^q(X)\bigr).
\]

\VerdierSemanticAnchor{ch4d-4-4-6-b}
\noindent\textit{b) 各項 $\mathrm{II}_{\infty}^{p,q}$ は、適切な
フィルトレーションを備えた対象 $F^{p+q}(X)$ に付随する次数付き対象である。}

\VerdierSemanticAnchor{ch4d-4-4-6-proof}
実際、(4.3.4.1) により標準同型
\[
X(p,p+1)\xrightarrow{\sim}H^p(X)[-p],
\]
がある。したがって（第 II 章、4.3.3.2）を用いると
\[
\mathrm{II}_2^{p,q}
=F^{p+q}\bigl(H^q(X)[-q]\bigr)
=F^p\bigl(H^q(X)\bigr),
\]
となり、a) の（記法を濫用した）等式が得られる。主張 b) は
（第 II 章、4.3.2）から従う。

\VerdierSemanticAnchor{ch4d-4-5}
\medskip\noindent\textbf{4.5. 収束の問題。}

\VerdierSemanticAnchor{ch4d-4-5-1}
\noindent\textbf{定義 4.5.1.}\quad
$\mathcal A$ と $\mathcal B$ を二つのアーベル圏とする。コホモロジー的関手
$F\colon D^*(\mathcal A)\longrightarrow\mathcal B$
（$*=+,-,b,{}$ ``空''）が\emph{右定常}（それぞれ\emph{左定常}）であるとは、
ある整数 $n_0$ が存在し、$D^*(\mathcal A)$ の各対象 $X$ について、
すべての $i<0$（それぞれ $>0$）に対し $H^i(X)=0$ ならば、すべての
$n<n_0$（それぞれ $n>n_0$）に対し $F^n(X)=0$ となることをいう。関手
$F$ が右定常かつ左定常ならば、\emph{定常}であるという。完全関手
$G\colon D^*(\mathcal A)\longrightarrow D^{*'}(\mathcal B)$ が
\emph{右定常}（それぞれ\emph{左定常}、それぞれ\emph{定常}）であるとは、
関手 $G$ と標準コホモロジー的関手
$D^{*'}(\mathcal B)\longrightarrow\mathcal B$ の合成として得られる
コホモロジー的関手が右定常（それぞれ左定常、それぞれ定常）であることをいう。
\endgroup
